\documentclass[11pt]{article}
\usepackage[a4paper,margin=24mm]{geometry}
\usepackage{amsmath,amssymb,amsthm}
\newtheorem{theorem}{Theorem}[section]
\newtheorem{remark}[theorem]{Remark}
\title{GST Direct Higher-Primitive Attack:
Algebraic Hodge Descent, Nonlinear Rees Logarithms,
and Genuine K-Symbol Cycle Birth}
\date{23 October 2026}
\begin{document}\maketitle

\section{Input and unconditionally constructed algebraic differential data}
Let X be ANY smooth projective complex d-fold,
let 2<=p<=d and let alpha be a PRESCRIBED rational
Hodge (p,p) class in H^(2p)(X,Q).
We accept the established GST Hodge-pure dimensional
fold theorem. The target remains a genuine
original-X rational CHOW cycle of codimension p
with canonical Betti class alpha.

Take any finite Zariski AFFINE cover U_i of X.
Their finite intersections are affine because
X is separated.

\begin{theorem}[Every prescribed rational Hodge class produces FINITE algebraic differential cocycle data on original X]
There exist finitely many genuinely ALGEBRAIC
forms
\[
 \boxed{
 \omega_\alpha=\{
 \omega^{a,b}_{i_0\dots i_a}\},\qquad
 a+b=2p,\quad b\ge p,\quad
 \omega^{a,b}_{i_0\dots i_a}
 \in\Gamma(U_{i_0\dots i_a},\Omega_X^b)
 }                                                     \tag{1}
\]
satisfying the actual Cech--de Rham
differential equation
\[
 \boxed{D\omega_\alpha
         =(\delta+(-1)^a d)\omega_\alpha=0.}      \tag{2}
\]
Their class under the canonical algebraic
de Rham-to-Betti comparison is alpha.
No codimension-p Chow cycle, algebraic
K0 class, or polynomial center for alpha
is inserted as an input.
\end{theorem}
\begin{proof}
For projective smooth X the actual algebraic
Hodge-to-de Rham spectral sequence degenerates
at E1. The image of the filtered hypercohomology
$\mathbb H^{2p}(X,\Omega_X^{\ge p})$ is
$F^pH^{2p}_{dR}(X)$.
The prescribed pure (p,p) alpha lies in F^p
after rational Betti comparison; choose
a filtered hypercohomology lift.
Since affine intersections have no higher
coherent cohomology, the finite Cech-de Rham
total complex computes the filtered algebraic
hypercohomology. A finite total cocycle gives
exactly the forms in (1) satisfying (2)
and maps to alpha.
This constructs true algebraic differential
source data from alpha, although these data
are not assumed to define an algebraic CYCLE.
\end{proof}

\section{SAME polynomial fold equations generate the algebraic log wedge}
On any original-X coordinate chart with an ACTUAL
smooth codimension-p algebraic fold center
$Z=V(z_1,\dots,z_p)$, the nonlinear
Rees deformation is the true algebraic scheme
\[
 R[t,u_1,\dots,u_p]/(t u_i-z_i).
\]
On the open chart where all log denominators
are invertible, its differential identities are
\[
 \boxed{dz_i=t\,du_i+u_i\,dt,\qquad
 d\log z_i=d\log t+d\log u_i.}                  \tag{3}
\]

\begin{theorem}[All-codimension EXACT logarithmic fold residue]
For EVERY p>=1, the very same Rees
fold equations yield the actual algebraic
p-form identity
\[
 \boxed{
 \begin{aligned}
 \bigwedge_{i=1}^p d\log z_i
  &=\bigwedge_{i=1}^p d\log u_i\\
  &\quad+d\log t\wedge
   \sum_{j=1}^p(-1)^{j-1}
     \bigwedge_{i\ne j}d\log u_i .
 \end{aligned}}                                         \tag{4}
\]
Consequently the logarithmic residue on
the fold boundary t=0, in this log chart, is
\[
 \boxed{\operatorname{Res}_{t=0}
   \bigwedge_{i=1}^p d\log z_i
    =\sum_{j=1}^p(-1)^{j-1}
       \bigwedge_{i\ne j}d\log u_i.}             \tag{5}
\]
The order in each omitted wedge is increasing.
\end{theorem}
\begin{proof}
Insert (3) in the exterior product.
All terms with two copies of dlog t vanish.
The no-dlog-t term is the full u-wedge.
A term with dlog t in position j incurs
the sign (-1)^(j-1) when moved to the front.
Its coefficient is the logarithmic residue.
This proof holds in the full exterior algebra
for every p, not only tested dimensions.
\end{proof}

\section{Actual cycle multiplication from algebraic local K-symbol residues}
For a genuine algebraic regular center
Z with local parameters z_1,...,z_p,
the Milnor K-symbol
\[
 \boxed{\{z_1,\dots,z_p\}\in
              K_p^M(\mathbb C(X))}
                                                               \tag{6}
\]
has logarithmic differential
$\bigwedge_i dz_i/z_i$.
The successive tame residues along the
regular ordered flag defined by the
z_i yield the unit multiplicity in
the terminal K0 at the generic point of Z.
This is the algebraic local intersection
multiplicity one, not a Hodge operator
being called a cycle.

\begin{theorem}[The polynomial fold creates a genuine Chow cycle from its K-residue center]
For any smooth complete-intersection center
Z of codimension p>=2, its true Milnor
symbol data and its normal Rees algebra
are built from the SAME regular equations
z_i=0. The residue has unit fundamental
multiplicity, and the actual blowup
$b:\widetilde X\to X$ with exceptional
divisor E constructs
\[
 \boxed{(-1)^{p-1}b_*[E]^p=[Z]
                             \in CH^p(X).}       \tag{7}
\]
The canonical Betti class of this genuine
original-X Chow cycle is the algebraic
fundamental class of Z.
\end{theorem}
\begin{proof}
The iterated local valuation residues of
an ordered regular-parameter symbol have
terminal value 1. The Chow fundamental
cycle of the regular complete intersection
therefore has generic multiplicity one.
For the Rees blowup, the exact normal-bundle
self-intersection formula gives
$E^p=(-1)^{p-1}j_*(\zeta^{p-1})$,
where $q:E=P_Z(N)\to Z$ and
$q_*\zeta^{p-1}=1$.
Push forward along $b\circ j=i\circ q$
to obtain (7) in ACTUAL Chow.
\end{proof}

\section{Adversarial direct alpha-to-polynomial-center construction}
We now attempt the requested arbitrary-X bridge:
the unconditionally constructed
algebraic differential cocycle (1)
must give genuinely new algebraic
local K-symbol residues (6), not merely
an algebraic differential form.
The first candidate puts an alpha-selected
algebraic form into the graph/zero section
of a FIXED algebraic vector bundle V
and Rees-folds its zero locus.

\begin{theorem}[Universal fixed-bundle zero-locus obstruction]
For EVERY smooth projective original X,
every algebraic rank-r vector bundle V and
ANY alpha-dependent algebraic section
$s_\alpha\in H^0(X,V)$ with regular
codimension-r zero scheme,
the TRUE Chow and canonical Betti classes obey
\[
 \boxed{
 [Z(s_\alpha)]_{CH}=c_r^{CH}(V),
 \qquad
 \mathrm{cl}^{can}_X[Z(s_\alpha)]
        =c_r^{Betti}(V),
 }                                                     \tag{8}
\]
INDEPENDENT of alpha.
For nonregular sections, the localized
top-Chern cycle pushed to X is likewise
$c_r(V)$. Thus the alpha-dependent
algebraic differential cocycle,
merely used as a variable section of
a fixed bundle followed by polynomial
Rees folding, does NOT generate new
Hodge-dependent native Chow directions.
\end{theorem}
\begin{proof}
The refined Gysin intersection of
the graph of a section of V and the
zero section computes the localized
Euler/top Chern class.
A regular section gives its
fundamental zero-scheme cycle; its
Chow class is c_r(V) cap [X].
For nonregular sections, proper
pushforward of the localized
top Chern class still equals
c_r(V) cap [X].
The bundle V is fixed,
so its Chow and Betti Chern
classes cannot vary with alpha.
This applies to the actual graph
and to its normal-cone Rees fold.
\end{proof}

\section{Exact outcome of this arbitrary-X route}
We have produced an ACTUAL finite algebraic
differential source from EVERY rational
Hodge alpha without a supplied Chow class;
we have translated the SAME nonlinear
dimension fold to an ALL-p algebraic
logarithmic residue and local K-symbol;
and we have proved (8), a direct
countertest to the tempting source
construction using only variable
forms in a fixed bundle.

The positive missing arrow is precisely
an INDEPENDENT construction from the
alpha-indexed filtered algebraic
differential cocycle to a NEW
algebraic vector bundle/regular ideal
whose logarithmic K-residue has
canonical Betti class alpha.
Equation (8) rules out one
specific fixed-bundle implementation,
not the GST geometric program.
Nothing here equates purity
of a Hodge decomposition with an
actual original-X Chow class
without such a source construction.
\end{document}